\scanpage{19}\mspage{18}
% @ first line
la \deleted{\unsure{raffiner en}} \unsure{renforcer}. \deleted{\illegible{1}}

% @ first, crossed condition
\deleted{\uline{\unsure{Contr|Cond}}\textsubscript{2} Si $X\in\unsure{F_{\grave e l}|F_{\acute e l}|F_{el}}$, alors l'ens \unsure{des} \unsure{$X'\ll X$}}
\[
\deleted{\{X'\in F_{el}\mid X'\ll \unsure{F_{el}|X}\}}
\]

% @ second condition
\uline{Contr}\textsubscript{2} Pour \unsure{tt} $X\in F_{\bar e l}$, considérons la dim comb. (pour $\leq$) des $X'\in F_{el}$ \deleted{tels que} qui raffinent $X$. Alors les dim. comb. de ces $X'$ sont majorées.

% @ dimension, middle
\unsure{Leur} borne sup. s'appelle la \uline{dimension \inserted{(\unsure{géom})} de $X$}. \deleted{\unsure{U}}La dim géom d'une figure, est par déf. le sup des dim géom. des \deleted{\unsure{ses}} figures élémentaires qui \unsure{la|les} constituent.
% slanted
\margin{\deleted{\unsure{Donc|D'où} la dimension}}

% @ interline chain
\inserted{C'est donc le sup des longueurs des chaînes $(X_0,X_1,\ldots,X_n)$ avec
\[
 X_0<X_1\ \ldots<X_{n-1}\ll X_n\in F.
\]}
% @ right edge below chain
\margin{\unsure{Elle|Ell} \deleted{$\unsure{F}$} peut être infinie pour $F$ \unsure{infinie}. Si $F'$ raffine $F$, $\unsure{\dim F'|\dim F}$ $\unsure{\leq|=}$ $\dim F$}

% @ dated paragraph
\margin{(\uline{11 juin})}
On aurait pu développer \unsure{un peu|un jour|une fois} les notions de ``partie constructible'', et adhérence, intérieur, frontière, composantes connexes associés à une figure, \unsure{avec} les \unsure{seules} \unsure{données|donnés} $\leq$, compatibilité sur \deleted{\unsure{la}} \deleted{\unsure{un}} $F_{\text{él.}}$.
% @ lower oblique note
% slanted
\margin{Aussi, la notion de \uline{structure de \unsure{contrée}} induite sur une \uline{\unsure{multistrate|multistr.}} figure. La notion de \unsure{sous-figure|sous figure} ($\unsure{\leq}$), \uline{composantes} \unsure{d'un} \unsure{univers|ensemble} \unsure{des} \unsure{points|parties} \unsure{pour} \unsure{``disj.''} \unsure{d'une} \unsure{figure}, \unsure{connexes|maximales}, celle \unsure{de|des} \illegible{1}, \unsure{Aussi}.}
% @ closing ordinary lines
J'y reviendrai sûrement. Pour le moment, j'ai envie de pousser \deleted{\unsure{plutôt}} un peu \unsure{ds|dans} la direction $\ll$, compatibilité,
